
\section{Guiding Principles and Revised Rules for the Lutheran Free Church}

\subsection{Guiding Principles}

1. The congregation is, according to God’s Word, the proper form of the Kingdom of God on earth.\footnote{When Principle No. 1 was discussed at the Free Church’s annual meeting in 1898, Prof. Sverdrup, according to Folkebladet for June 15, 1898, spoke as follows:

It is of great importance that one declare oneself on these matters; for here it is not a matter of adopting points, but of our also acknowledging from the heart what we confess in these rules. It is something new when in Point 1 we declare the congregation to be the proper form of the Kingdom of God on earth; for for hundreds of years men have called the Church the proper form.

If, however, we adopt this point, then we go back to the time of the Apostles and to God’s Word.

According to God’s Word, the congregation is not a society governed by Pope, King, Council, or any other authority; but the congregation is the gathering of believing people around God’s Word and the means of grace. There is not in God’s Word a single sign that any administration should be set over the congregation to rule over it.}

2. The congregation consists of believing people who, by the use of the means of grace and the gifts of grace according to the direction of God’s Word, seek their own salvation and eternal blessedness and that of all their fellow men.

3. According to the New Testament, the congregation requires an outward organization with a register of its members, election of office-bearers, fixed times and places for its gatherings, etc.

4. Since not all members of the congregation’s outward organization are always believing people, and since such sinners often take false comfort from their outward connection with the congregation, it is the congregation’s holy calling to purify itself both by awakening proclamation of God’s Word, by earnest admonition and reminder, and by exclusion of manifest and obstinate sinners.\footnote{Point 4 in the Principles originally read as follows:
“Ordinarily the congregation’s outward organization also includes people who in God’s Word are called dead and sleeping, wherefore within the congregation too the cry of awakening and the preaching of repentance are needed, just as it is the believers’ holy calling also within the congregation to be light and salt, lest the outward connection with the congregation become a false comfort and a hindrance.

Since there were some who found a contradiction between Points 2 and 4, Point 4 was altered after Sverdrup had, according to Folkebladet’s report for June 16, 1897, made the following statement:
As regards the contradiction between the two points 2 and 4, it is the same opposition as in every Christian—the law of the Spirit and the law of the flesh. The congregation consists of those who belong to the body of Jesus Christ—that which lives—what else should it consist of? But the outward organization has dead members in it. God’s Word says this too. There was a Judas even among the twelve. The works of the flesh are knotty things to have to do with. Those who try to grapple fast with them injure themselves most and the works of the flesh least. I should wish there were a Word of God that gave another view of the congregation than this Point 4, but I cannot find any. There are vessels for honor and vessels for dishonor in every house. The congregation consists and must consist of living members; but there are dead ones also, and these we cannot by our understanding and our sight separate out. But the congregation does not consist of dead ones. There are many unclean vessels and creatures in a house, but the house does not consist of these unclean things. Can anyone then, for example, sit in judgment upon the children in the congregation, and say, when the children fall out of their baptismal covenant, that they no longer belong to the congregation? Such a congregation I should at any rate leave as soon as possible.—Thus stand the works of the flesh.

In his lectures Sverdrup often used this as a definition of the word “consist”: “that which really holds up, even though much else hangs on.”}


5. The congregation governs its own affairs under the authority of God’s Word and Spirit, and acknowledges no other ecclesiastical authority or power over itself.

6. The free congregation esteems and loves all the spiritual gifts which the Lord grants for its edification, and seeks to enliven the gifts and promote their use.

7. The free congregation receives with gladness the mutual help which the congregations can render one another in the work of advancing the Kingdom of God.

8. Such mutual help consists partly in this, that the spiritual powers of one congregation come the other to good through free meetings, mutual visits, lay activity, etc.; partly in this, that the congregations voluntarily and according to the prompting of God’s Spirit work together for the solution of such tasks as may exceed the powers of a single congregation.

9. Among such purposes are named in particular a seminary for pastors, the dissemination of Bibles and other writings and papers, Home Mission, mission to the heathen, mission to the Jews, deaconess homes, orphanages, and other works of mercy.

10. Free congregations have no right to require that other congregations submit to their opinion, will, judgment, or determination, wherefore every dominion of a majority of congregations over a minority is to be rejected.

11. The common organs which may be found desirable for the cooperation of the congregations, such as larger or smaller meetings, committees, office-bearers, etc., cannot in a Lutheran Free Church impose upon the single congregation any duty, bond, compulsion, or burden, but are only authorized to come forward with proposals and exhortations to congregations and persons.

12. Every free congregation has, like each individual believer, under the impulse of God’s Spirit and by the right of love, the right to do good and to labor for the salvation of souls and the awakening of spiritual life, as far as its ability and strength reach; in such free, spiritual activity it is restricted neither by parish ties nor by denominational boundaries.


\pagebreak
\subsection{The Revised Rules for Work}

1. The name of this association is: The Lutheran Free Church.\footnote{Originally the ‘Rules for Work’ read as follows:

1. The name of this association is: The Lutheran Free Church.

2. Its purpose is to labor for the enlivening and liberation of all Norwegian Lutheran congregations, so that, according to their calling and their ability, they may work in spiritual independence and freedom for the cause of God’s Kingdom at home and abroad through such organs and institutions as the congregations themselves might determine.

3. This purpose is sought attained chiefly by the holding of greater and smaller meetings, by the dissemination of suitable literature, by the organization of free-church committees and associations, by the sending out of emissaries, and by other means which from time to time might prove suitable.

4. The Free Church consists of Lutheran congregations which, by resolution, acknowledge its principles and rules and report this to the Secretary, regardless of other denominational connection.

5. The Free Church holds an annual meeting, which as a rule begins on the first Wednesday in the month of June, and which elects the necessary committees, emissaries, and officers, and determines which ecclesial matters the meeting will particularly recommend to the congregations.

6. At the Free Church’s annual meeting, all voting members of the congregations that constitute the Free Church have the right to vote, insofar as they appear from those congregations; likewise members entitled to vote of other Lutheran congregations who, by submitting their names to the Secretary of the annual meeting on forms specially prepared for that purpose, declare that they approve the principles and rules of the Free Church and will labor for the purpose stated for it in §2.

7. At the Free Church’s annual meeting, reports are heard from officers and committees, as well as from such ecclesial institutions as the annual meeting grants access thereto.

8. The officers of the Free Church shall be a Chairman and a Secretary. They are elected by the annual meeting for a term of one year. The Chairman is the presiding officer of the annual meeting; the Secretary keeps the minutes of the annual meeting; together these two officers issue the annual report of the Free Church, determine the time and place of the annual meeting, unless the immediately preceding annual meeting has determined it, and see to it that the annual meeting is announced at least two months before it is held.

9. Each year the annual meeting elects an ordinator, who shall assist the congregations by ordaining theological candidates who have been duly called to be pastors; likewise, each year it elects a committee for Inner Mission, and a committee for mission to the heathen, and a church organizational committee.

10. Besides the regular annual meeting, during the course of the month of June each year local free-church meetings are held in the various congregations, which may invite thereto. At each of these meetings, if possible, one or more of the Free Church’s functionaries or sent workers shall be present; and particular weight shall be laid upon conversations and deliberations of an edifying character, intended to cast light upon the work of solving the free-church tasks, and especially upon the enlivening of the congregation among us.

11. These rules may be altered in such a way that a proposal for alteration is brought forward at one annual meeting and voted upon at the next annual meeting; its adoption then requires a two-thirds majority of the votes cast.
}



% The section starts here!



2. Its purpose is to labor for the vivification and liberation of all Norwegian Lutheran congregations, so that according to their calling and their ability they may labor in spiritual independence and freedom for the cause of the Kingdom of God at home and abroad through such organs and institutions as the congregations themselves may determine.

3. This purpose is sought attained chiefly by the training of men and women for spiritual work in and for the congregation, by the holding of larger and smaller meetings, by the dissemination of suitable literature, by the organization of voluntary committees and associations, by the sending out of emissaries, and by other means which may from time to time prove suitable.

4. The Free Church consists of congregational members who in its church order unreservedly assent to the ancient ecumenical Creeds, Luther’s Small Catechism, and the Unaltered Augsburg Confession, and who by formal decision acknowledge the Free Church’s principles and rules and report this to the Secretary, notwithstanding other ecclesial affiliation.

5. The Free Church holds an annual meeting, which as a rule begins on the second Wednesday of the month of June, and which elects the necessary committees, emissaries, and officers, and determines which ecclesiastical tasks the meeting will particularly recommend to the congregations.

6. The right to vote at the Free Church’s annual meetings belongs to all enfranchised congregational members who appear from the congregations that constitute the Free Church; likewise, enfranchised members of other Lutheran congregations who, by submitting their names to the Secretary of the annual meeting on forms specially prepared for that purpose, signify that they approve the Free Church’s principles and rules and will work for the purpose set forth for it in paragraph 2.

7. The officers of the Free Church shall be a Chairman, Vice-Chairman, and Secretary. They are elected by the annual meeting for a term of one year. The Chairman is the presiding officer of the annual meeting, the Secretary keeps the minutes of the annual meeting; these two officers together publish the Free Church’s annual report and determine the time and place of the annual meeting, unless the immediately preceding annual meeting has determined this, and also see to it that the annual meeting is publicly announced at least two months before it is held.

8. The annual meeting elects an Organizational Committee of three members. One member is elected each year for a term of three years.

9. a) The organization committee labors to make the Free Church’s principles and rules known and discussed in the congregations, so that the Free Church’s task may gradually come to be more clearly understood and better carried through.

b) It labors that the congregations may be as numerously and generally represented at the Free Church’s annual meetings as possible.

c) It provides for the holding of larger meetings, and arranges the time, place, program, and personnel for the same.

d) The committee seeks to bring about a general interchange of the spiritual powers within the congregations.

e) It will assist congregations and pastors with counsel and guidance where this is desired.

f) Each year it gathers and reports to the annual meetings as complete information as possible concerning the Free Church’s work in its various branches, such as congregational work, foreign and home mission, Sunday schools, women’s societies, youth societies, church buildings, church dedications, etc.

10. The annual meeting elects each year an ordinator, who shall assist the congregations in ordaining theological candidates who in due manner have been called as pastors.

11. The annual meeting elects a mission superintendent, who shall devote all his time to the service of Home Mission and assist the Board of Home Missions in organizing the work of Home Mission. He is elected for three years, and his salary is determined by the Board of Home Missions.

12. The annual meeting elects a school board for Augsburg Seminary and nominates members to the Lutheran Board of Missions, Board of Home Missions, the Corporation of Augsburg Seminary, and its Board of Trustees.

13. Chairman, Secretary, Ordinator, and chairmen of standing committees cannot be re-elected for more than two consecutive terms of office.

14, At the Free Church’s annual meetings reports are heard from the organization committee, the Lutheran Board of Missions, the Board of Home Missions, and Augsburg Seminary, and from committees appointed by the Free Church’s annual meeting, as well as from such ecclesiastical institutions as the annual meeting gives admission thereto.

15. These rules may be changed in such manner that a proposal for change is introduced at one annual meeting and voted upon at the next annual meeting; its adoption then requires a two-thirds majority of the votes cast.


